\chapter{{\fontsize{14pt}{21pt}\selectfont\MakeUppercase{Практическая часть и технологии}}}
\label{ch:technologies}

\section{Технологии и инструменты}

В~рамках проекта используются следующие языки программирования, фреймворки и инфраструктурные компоненты.

\textbf{Серверная часть и агенты:} язык \textbf{Go} (1.24+), HTTP-маршрутизатор go-chi/chi, очереди сообщений \textbf{NATS}, протокол взаимодействия агентов \textbf{A2A} (Agent-to-Agent) поверх NATS (request-reply). Для платформ без публичного API применяется браузерная автоматизация \textbf{Playwright} (Go-биндинг \texttt{playwright-go}).

\textbf{Клиентская часть:} \textbf{Next.js} 14 на базе React~18, \textbf{Tailwind~CSS} и \textbf{shadcn/ui}. Выбор обоснован поддержкой SSR, интеграцией с~Server-Sent Events для потоковых ответов оркестратора и обширной экосистемой компонентов.

\textbf{Данные:} \textbf{PostgreSQL} 16 (реляционные данные: пользователи, бизнес-объекты, интеграции, подписки), \textbf{MongoDB} 7 (гибкая схема: диалоги, сообщения, задачи, отзывы, публикации), \textbf{Redis} 7 (сессии, лимиты, кэш), \textbf{MinIO} (S3-совместимое хранилище файлов).

\textbf{Инфраструктура:} Docker Compose (локальная разработка), Kubernetes (промышленное развёртывание), GitHub Actions (CI/CD). Мониторинг: Prometheus, Grafana; логирование: структурированные логи (slog), при необходимости Loki.

\section{Протокол взаимодействия агентов (A2A)}

Все агенты реализуют единый интерфейс и общаются по протоколу A2A: оркестратор отправляет \texttt{ToolRequest} (идентификатор задачи, имя инструмента, аргументы, идентификатор бизнеса), агент выполняет операцию (через API платформы или RPA) и возвращает \texttt{ToolResponse} (успех/ошибка, результат). Транспорт~--- NATS, подписки вида \texttt{tasks.telegram}, \texttt{tasks.vk}, \texttt{tasks.yandex\_business}. Метод интеграции (REST API или Playwright) инкапсулирован внутри агента и прозрачен для оркестратора.

\section{Спецификации агентов (MVP)}

\textbf{VK (API):} официальный VK API, OAuth~2.0; операции: публикация постов в~сообщество, обновление информации сообщества, чтение и ответ на комментарии. Инструменты: \texttt{vk\_\_publish\_post}, \texttt{vk\_\_update\_group\_info} и~др.

\textbf{Telegram (API):} Telegram Bot API, токен бота; операции: публикация в~канал, редактирование и закрепление сообщений, уведомления владельцу. Инструменты: \texttt{telegram\_\_send\_channel\_post}, \texttt{telegram\_\_send\_notification} и~др.

\textbf{Яндекс.Бизнес (RPA):} публичного API нет, используется браузерная автоматизация Playwright (безголовый Chromium). Операции: обновление часов работы и контактов, загрузка фото, чтение отзывов и ответ на них. Требования: устойчивые селекторы, задержки 1--5~с между действиями, снимок экрана при ошибке, повтор с экспоненциальной задержкой.

Выбор языков и фреймворков обоснован методом взвешенных критериев (производительность и конкурентность, скорость прототипирования, экосистема, развёртывание, сопровождаемость); Go и Next.js набрали наивысшие баллы в~своих категориях~\cite{src:17}.
